\chapter{Curvilinear Coordinate Systems \& Scale Factors}
\index{Coordinate Systems}
\index{Cartesian Coordinates}
\index{Cylindrical Coordinates}
\index{Spherical Coordinates}
\index{Scale Factors}
\index{Differential Elements}
\index{Metric Tensor}
\index{Lamé Coefficients}

% ==========================================================
% Chapter 1 Instructional Management Plan & Syllabus
% ==========================================================
\section*{Chapter 1 Instructional Management Plan}
\addcontentsline{toc}{section}{Chapter 1 Instructional Management Plan}

\begin{planbox}[colframe=rbruNavy, colbacktitle=rbruNavy]{1. Chapter Topics and Core Content}
\begin{enumerate}[topsep=2pt, itemsep=1pt, partopsep=0pt, parsep=0pt, leftmargin=1.5em]
    \item Introduction and Cartesian Coordinate Foundations
    \item Cylindrical Coordinate System
    \item Spherical Coordinate System and GPS Geodesic Applications
    \item Summary of Key Formulas and Coordinate Transformations
    \item Chapter Exercises and Worked Physical Problems
\end{enumerate}
\end{planbox}

\begin{planbox}[colframe=rbruGreen, colbacktitle=rbruGreen]{2. Behavioral Learning Objectives (OBE)}
\begin{enumerate}[topsep=2pt, itemsep=1pt, partopsep=0pt, parsep=0pt, leftmargin=1.5em]
    \item Rigorously define and transform between Cartesian, cylindrical, and spherical coordinates.
    \item Derive metric scale factors (Lamé coefficients $h_i$) for orthogonal curvilinear systems.
    \item Formulate differential displacement, surface, and volume elements across coordinate bases.
    \item Apply spherical coordinates to satellite geodesics and GPS positioning geometry.
\end{enumerate}
\end{planbox}

\newpage

\begin{planbox}[colframe=rbruNavy!85!black, colbacktitle=rbruNavy!85!black]{3. Teaching Methods and Instructional Activities}
\begin{enumerate}[topsep=2pt, itemsep=1pt, partopsep=0pt, parsep=0pt, leftmargin=1.5em]
    \item Theoretical lecture integrated with 3D dynamic coordinate plane visualizations.
    \item Collaborative problem-solving sessions on metric tensor and scale factor derivations.
    \item Computational laboratory implementing coordinate transformations and plotting with Python.
    \item Case study discussion on GPS satellite positioning and orbital geometry.
\end{enumerate}
\end{planbox}

\begin{planbox}[colframe=rbruGold!90!black, colbacktitle=rbruGold!90!black]{4. Instructional Media and Educational Technology}
\begin{enumerate}[topsep=2pt, itemsep=1pt, partopsep=0pt, parsep=0pt, leftmargin=1.5em]
    \item Course textbook: Chapter 1 Curvilinear Coordinate Systems \& Scale Factors.
    \item Electronic lecture slides with 3D interactive vector models.
    \item Python scientific simulation scripts using NumPy and Matplotlib.
    \item Digital learning management system and online computational exercises.
\end{enumerate}
\end{planbox}

\begin{planbox}[colframe=rbruSlate, colbacktitle=rbruSlate]{5. Measurement and Learning Assessment}
\begin{enumerate}[topsep=2pt, itemsep=1pt, partopsep=0pt, parsep=0pt, leftmargin=1.5em]
    \item Formative evaluation of in-class engagement and active problem-solving participation.
    \item Assessment of homework problem sets on coordinate transformations and metric elements.
    \item Chapter 1 diagnostic quiz evaluating analytical and conceptual mastery.
\end{enumerate}
\end{planbox}
\clearpage

\section{Introduction and Cartesian Coordinate Foundations}
In theoretical physics, formulating dynamic laws, field distributions, and conservation principles requires a rigorous frame of reference. Selecting a coordinate system whose geometric symmetry reflects that of the physical boundary conditions significantly reduces the analytical complexity of governing differential equations.

The Cartesian coordinate system $(x, y, z)$ represents the fundamental baseline. Three mutually orthogonal axes intersect at the origin $O(0,0,0)$, forming a right-handed basis characterized by fixed unit vectors $\hat{i}, \hat{j}, \hat{k}$ satisfying the orthonormality condition:
\begin{equation}
\hat{e}_i \cdot \hat{e}_j = \delta_{ij} = \begin{cases} 1, & i = j \\ 0, & i \neq j \end{cases}
\end{equation}
where $\delta_{ij}$ is the Kronecker delta. The position vector $\vec{r}$ of a point $P(x, y, z)$ is:
\begin{equation}
\vec{r} = x\hat{i} + y\hat{j} + z\hat{k}
\end{equation}
The differential displacement vector and squared infinitesimal arc length are:
\begin{align}
d\vec{r} &= dx\,\hat{i} + dy\,\hat{j} + dz\,\hat{k} \\
ds^2 &= |d\vec{r}|^2 = dx^2 + dy^2 + dz^2
\end{align}
The differential volume element is simply $dV = dx\,dy\,dz$.

\begin{mathdefinition}[Orthogonal Curvilinear Coordinates and Scale Factors]
Let $(u_1, u_2, u_3)$ be generalized coordinates related to Cartesian coordinates by invertibly differentiable functions $x = x(u_1, u_2, u_3)$, $y = y(u_1, u_2, u_3)$, and $z = z(u_1, u_2, u_3)$. The system is \textbf{orthogonal} if the local basis vectors are mutually perpendicular. The \textbf{metric scale factors} (Lamé coefficients) are defined as:
\begin{equation}
h_i = \left|\frac{\partial \vec{r}}{\partial u_i}\right| = \sqrt{\left(\frac{\partial x}{\partial u_i}\right)^2 + \left(\frac{\partial y}{\partial u_i}\right)^2 + \left(\frac{\partial z}{\partial u_i}\right)^2}
\end{equation}
The unit basis vectors are $\hat{e}_i = \frac{1}{h_i}\frac{\partial \vec{r}}{\partial u_i}$. The differential line and volume elements become:
\begin{align}
ds^2 &= h_1^2 du_1^2 + h_2^2 du_2^2 + h_3^2 du_3^2 \\
dV &= h_1 h_2 h_3\,du_1\,du_2\,du_3
\end{align}
For Cartesian coordinates, $h_x = h_y = h_z = 1$.
\end{mathdefinition}

\section{Cylindrical Coordinate System}
Physical configurations exhibiting axial or rotational symmetry around a preferred axis---such as magnetic fields produced by long solenoids, laminar fluid flow through pipes, or rotating flywheels---are most naturally represented in cylindrical coordinates $(\rho, \phi, z)$.

Here, $\rho \in [0, \infty)$ is the perpendicular radial distance from the $z$-axis, $\phi \in [0, 2\pi)$ is the azimuthal angle measured counterclockwise from the positive $x$-axis, and $z \in (-\infty, \infty)$ is the axial coordinate:
\begin{equation}
x = \rho\cos\phi, \quad y = \rho\sin\phi, \quad z = z
\end{equation}
Inverting these relations yields:
\begin{equation}
\rho = \sqrt{x^2 + y^2}, \quad \phi = \arctan\left(\frac{y}{x}\right), \quad z = z
\end{equation}

To compute the scale factors, we differentiate $\vec{r}(\rho, \phi, z) = \rho\cos\phi\,\hat{i} + \rho\sin\phi\,\hat{j} + z\,\hat{k}$:
\begin{align}
\frac{\partial \vec{r}}{\partial \rho} &= \cos\phi\,\hat{i} + \sin\phi\,\hat{j} \implies h_\rho = 1 \\
\frac{\partial \vec{r}}{\partial \phi} &= -\rho\sin\phi\,\hat{i} + \rho\cos\phi\,\hat{j} \implies h_\phi = \rho \\
\frac{\partial \vec{r}}{\partial z} &= \hat{k} \implies h_z = 1
\end{align}
Thus, the orthonormal unit vectors are:
\begin{align}
\hat{\rho} &= \cos\phi\,\hat{i} + \sin\phi\,\hat{j} \\
\hat{\phi} &= -\sin\phi\,\hat{i} + \cos\phi\,\hat{j} \\
\hat{k} &= \hat{k}
\end{align}
The differential displacement, line element squared, and volume element are:
\begin{align}
d\vec{r} &= d\rho\,\hat{\rho} + \rho\,d\phi\,\hat{\phi} + dz\,\hat{k} \\
ds^2 &= d\rho^2 + \rho^2 d\phi^2 + dz^2 \\
dV &= \rho\,d\rho\,d\phi\,dz
\end{align}

\begin{pitfallbox}[Time Derivatives of Non-Cartesian Unit Vectors]
Unlike Cartesian unit vectors $(\hat{i}, \hat{j}, \hat{k})$ which remain invariant in space and time, cylindrical and spherical unit vectors depend on position. When differentiating with respect to time $t$:
\begin{equation}
\dot{\hat{\rho}} = \frac{d\hat{\rho}}{dt} = \dot{\phi}\,\hat{\phi}, \quad \dot{\hat{\phi}} = -\dot{\phi}\,\hat{\rho}
\end{equation}
Neglecting these geometric derivative terms produces severe errors in velocity and acceleration derivations.
\end{pitfallbox}

\section{Spherical Coordinate System}
For central force fields (e.g., Newtonian gravity, Coulomb potential) and spherically symmetric wave propagation, spherical coordinates $(r, \theta, \phi)$ are indispensable.

The coordinates are defined as:
\begin{itemize}
    \item $r \in [0, \infty)$: radial distance from the origin.
    \item $\theta \in [0, \pi]$: polar (zenith) angle measured from the positive $z$-axis.
    \item $\phi \in [0, 2\pi)$: azimuthal angle in the $xy$-plane from the positive $x$-axis.
\end{itemize}
The Cartesian transformation equations are:
\begin{equation}
x = r\sin\theta\cos\phi, \quad y = r\sin\theta\sin\phi, \quad z = r\cos\theta
\end{equation}
The metric scale factors are obtained from $\vec{r}(r, \theta, \phi)$:
\begin{align}
h_r &= \left|\frac{\partial \vec{r}}{\partial r}\right| = 1 \\
h_\theta &= \left|\frac{\partial \vec{r}}{\partial \theta}\right| = r \\
h_\phi &= \left|\frac{\partial \vec{r}}{\partial \phi}\right| = r\sin\theta
\end{align}
Consequently, the differential elements are:
\begin{align}
d\vec{r} &= dr\,\hat{r} + r\,d\theta\,\hat{\theta} + r\sin\theta\,d\phi\,\hat{\phi} \\
ds^2 &= dr^2 + r^2 d\theta^2 + r^2\sin^2\theta\,d\phi^2 \\
dV &= r^2\sin\theta\,dr\,d\theta\,d\phi
\end{align}

\begin{table}[H]
\centering
\small
\begin{tabularx}{\textwidth}{lcccc}
\toprule
\textbf{Coordinate System} & \textbf{Coordinates} & \textbf{Scale Factors} $(h_1, h_2, h_3)$ & \textbf{Infinitesimal Volume} $dV$ \\
\midrule
Cartesian & $(x, y, z)$ & $(1, 1, 1)$ & $dx\,dy\,dz$ \\
Cylindrical & $(\rho, \phi, z)$ & $(1, \rho, 1)$ & $\rho\,d\rho\,d\phi\,dz$ \\
Spherical & $(r, \theta, \phi)$ & $(1, r, r\sin\theta)$ & $r^2\sin\theta\,dr\,d\theta\,d\phi$ \\
\bottomrule
\end{tabularx}
\caption{Comparison of metric scale factors and volume elements across standard coordinate systems.}
\label{tab:ch01_scale_factors}
\end{table}

\section{Worked Examples and Physical Applications}

\begin{psctexample}[Kinematic Acceleration in Cylindrical Coordinates]
A point particle moves in a plane ($z = 0$) such that its radial distance is $\rho(t) = \rho_0 e^{\alpha t}$ while rotating with a constant angular speed $\omega = \dot{\phi}$. Determine the general acceleration vector $\vec{a}(t)$ in cylindrical basis.

\textbf{\color{rbruNavy}Picture / Conceptualization:}
The particle spirals outward in the $xy$-plane. The coordinate basis vectors $\hat{\rho}$ and $\hat{\phi}$ rotate with the particle at rate $\omega$.

\textbf{\color{rbruNavy}Strategy / Governing Laws:}
Express position $\vec{r} = \rho\hat{\rho}$. Differentiate twice with respect to time, applying the kinematic chain rule to account for $\dot{\hat{\rho}} = \dot{\phi}\hat{\phi}$ and $\dot{\hat{\phi}} = -\dot{\phi}\hat{\rho}$.

\textbf{\color{rbruNavy}Calculation / Analytical Solution:}
1. Velocity vector:
\begin{equation}
\vec{v} = \dot{\vec{r}} = \dot{\rho}\hat{\rho} + \rho\dot{\hat{\rho}} = \dot{\rho}\hat{\rho} + \rho\dot{\phi}\hat{\phi}
\end{equation}
2. Acceleration vector:
\begin{align}
\vec{a} &= \frac{d\vec{v}}{dt} = \ddot{\rho}\hat{\rho} + \dot{\rho}\dot{\hat{\rho}} + (\dot{\rho}\dot{\phi} + \rho\ddot{\phi})\hat{\phi} + \rho\dot{\phi}\dot{\hat{\phi}} \nonumber\\
&= (\ddot{\rho} - \rho\dot{\phi}^2)\hat{\rho} + (2\dot{\rho}\dot{\phi} + \rho\ddot{\phi})\hat{\phi}
\end{align}
3. Substituting $\rho(t) = \rho_0 e^{\alpha t}$, $\dot{\rho} = \alpha\rho$, $\ddot{\rho} = \alpha^2\rho$, $\dot{\phi} = \omega$, and $\ddot{\phi} = 0$:
\begin{equation}
\vec{a}(t) = (\alpha^2 - \omega^2)\rho_0 e^{\alpha t}\hat{\rho} + 2\alpha\omega\rho_0 e^{\alpha t}\hat{\phi}
\end{equation}

\textbf{\color{rbruNavy}Think / Physical Insights:}
The radial component contains the centrifugal acceleration term $-\rho\omega^2$, while the azimuthal component $2\alpha\omega\rho$ is the Coriolis acceleration. When $\alpha = \omega$, the radial acceleration vanishes, yielding purely tangential acceleration.
\qedexam
\end{psctexample}

\begin{pythonbox}[Trajectory and Scale Factor Verification in Python]
import numpy as np
import matplotlib.pyplot as plt

# Parameters for spiral motion
t = np.linspace(0, 2, 500)
rho_0 = 1.0
alpha = 0.8
omega = 4.0

rho = rho_0 * np.exp(alpha * t)
phi = omega * t

x = rho * np.cos(phi)
y = rho * np.sin(phi)

# Acceleration components
a_rho = (alpha**2 - omega**2) * rho
a_phi = 2 * alpha * omega * rho

print(f"At t = 1.0s: rho = {rho_0*np.exp(alpha):.3f} m")
print(f"Radial acceleration: {a_rho[250]:.3f} m/s^2, Tangential: {a_phi[250]:.3f} m/s^2")
\end{pythonbox}

\begin{psctexample}[GPS Geodesic Distance from Spherical Coordinates]
Two points on Earth's surface are given by geographic coordinates: Bangkok $(\lambda_1 = 13.75°\text{N},\; \phi_1 = 100.52°\text{E})$ and Tokyo $(\lambda_2 = 35.68°\text{N},\; \phi_2 = 139.65°\text{E})$. Using the spherical law of cosines, calculate the great-circle (geodesic) distance between them. Take Earth's mean radius $R_\oplus = 6371\;\text{km}$.

\textbf{Solution:}

\textbf{Step 1 — Convert to spherical coordinates.}
In physics conventions, the polar angle $\theta$ is measured from the North Pole, so $\theta = 90° - \lambda$:
\begin{align}
\theta_1 &= 90° - 13.75° = 76.25°, \quad \phi_1 = 100.52° \\
\theta_2 &= 90° - 35.68° = 54.32°, \quad \phi_2 = 139.65°
\end{align}

\textbf{Step 2 — Apply the spherical law of cosines.}
The angle $\Delta\sigma$ subtended at Earth's centre between two points on a unit sphere satisfies:
\begin{equation}
\cos(\Delta\sigma) = \cos\theta_1\cos\theta_2 + \sin\theta_1\sin\theta_2\cos(\phi_2 - \phi_1)
\end{equation}
Substituting $\Delta\phi = 139.65° - 100.52° = 39.13°$:
\begin{align}
\cos(\Delta\sigma) &= \cos(76.25°)\cos(54.32°) + \sin(76.25°)\sin(54.32°)\cos(39.13°) \\
                   &= (0.2385)(0.5835) + (0.9712)(0.8122)(0.7753) \\
                   &= 0.1391 + 0.6124 = 0.7515
\end{align}
Therefore $\Delta\sigma = \arccos(0.7515) = 41.27°= 0.7203\;\text{rad}$.

\textbf{Step 3 — Compute the geodesic distance:}
\begin{equation}
\boxed{d = R_\oplus \cdot \Delta\sigma = 6371 \times 0.7203 \approx 4589\;\text{km}}
\end{equation}

\textbf{\color{rbruNavy}Think / Physical Insights:}
This calculation underlies all GPS receiver algorithms. The Haversine formula (used in navigation software) is numerically superior to the simple cosine formula when $\Delta\sigma \to 0$, because it avoids catastrophic cancellation. For intercontinental distances, however, both formulas agree to better than $0.1\%$.
\qedexam
\end{psctexample}

\begin{pythonbox}[Great-Circle Distance Verification in Python]
import numpy as np

def great_circle(lat1, lon1, lat2, lon2, R=6371):
    # Convert degrees to radians
    t1, t2 = np.radians(90 - lat1), np.radians(90 - lat2)
    p1, p2 = np.radians(lon1), np.radians(lon2)
    cos_dsigma = np.cos(t1)*np.cos(t2) + np.sin(t1)*np.sin(t2)*np.cos(p2 - p1)
    dsigma = np.arccos(np.clip(cos_dsigma, -1, 1))
    return R * dsigma

d = great_circle(13.75, 100.52, 35.68, 139.65)
print(f"Bangkok → Tokyo great-circle distance: {d:.1f} km")
print(f"Subtended angle: {np.degrees(d/6371):.2f} degrees")
\end{pythonbox}

\section{Summary of Key Formulas}
\begin{table}[H]
\centering
\small
\begin{tabularx}{\textwidth}{llX}
\toprule
\textbf{Concept} & \textbf{Formula} & \textbf{Physical Significance} \\
\midrule
Metric Scale Factor & $h_i = |\partial\vec{r}/\partial u_i|$ & Local arc length per unit coordinate variation \\
Infinitesimal Line Element & $ds^2 = \sum_{i} h_i^2 du_i^2$ & Invariant metric distance in differential geometry \\
Cylindrical Basis Derivatives & $\dot{\hat{\rho}} = \dot{\phi}\hat{\phi}, \; \dot{\hat{\phi}} = -\dot{\phi}\hat{\rho}$ & Origin of centripetal and Coriolis inertial terms \\
Spherical Solid Angle & $d\Omega = \sin\theta\,d\theta\,d\phi$ & Measure of angular aperture in radiation and scattering \\
Great-Circle Distance & $d = R\arccos(\cos\theta_1\cos\theta_2 + \sin\theta_1\sin\theta_2\cos\Delta\phi)$ & Shortest path on a sphere (GPS, aviation) \\
\bottomrule
\end{tabularx}
\caption{Summary of coordinate geometry relations and physical interpretations.}
\label{tab:ch01_summary}
\end{table}

\section{Chapter Exercises}
\begin{enumerate}
    \item \textbf{Analytical Transformation:} Find the spherical coordinates $(r, \theta, \phi)$ of the point given in Cartesian coordinates by $P(-2, 2\sqrt{3}, -4)$.
    \item \textbf{Scale Factor Derivation:} Consider parabolic cylindrical coordinates $(u, v, z)$ defined by $x = \frac{1}{2}(u^2 - v^2)$, $y = uv$, $z = z$. Calculate scale factors $h_u, h_v, h_z$ and show that the coordinate system is orthogonal.
    \item \textbf{Solid Angle of a Cone:} A right-circular cone has a semi-vertical angle $\alpha$. By integrating the spherical solid angle $d\Omega = \sin\theta\,d\theta\,d\phi$, prove that the solid angle subtended at the apex is $\Omega = 2\pi(1 - \cos\alpha)$.
    \item \textbf{Velocity in Spherical Coordinates:} A particle moves on a sphere of radius $R$ with $\dot\theta = \omega_\theta$ (const) and $\dot\phi = \omega_\phi$ (const). Write the velocity vector $\vec{v}$ in the spherical basis $(\hat{r}, \hat{\theta}, \hat{\phi})$ and compute its magnitude. Under what condition is $|\vec{v}|$ constant?
    \item \textbf{Ellipsoidal Coordinates:} Prolate spheroidal coordinates $(\xi, \eta, \phi)$ are related to Cartesian coordinates by $x = a\sinh\xi\sin\eta\cos\phi$, $y = a\sinh\xi\sin\eta\sin\phi$, $z = a\cosh\xi\cos\eta$ where $a > 0$. Derive the three scale factors and the volume element $dV$.
\end{enumerate}
